\documentclass[11pt,a4paper]{article}
\usepackage[utf8]{utf8}
\usepackage[margin=1in]{geometry}
\usepackage{amsmath,amssymb}
\usepackage{enumitem}

\begin{document}

\begin{center}
    {\Large \textbf{U.G. Semester II Examination 2024--25 (NEP)}}\\[4pt]
    {\large \textbf{SOCIOLOGY}}\\[4pt]
    \textbf{Paper Code: SOCMD21 : Indian Knowledge System}\\[6pt]
    \textbf{Time: 2:30 Hours} \hfill \textbf{Full Marks: 70}
\end{center}
\hr

\noindent \textbf{Note:} Write your Roll No. at the top immediately on the receipt of this question paper. The paper consists of three sections: A, B, and C.

\section*{SECTION -- A (Objective Type Questions)}
\noindent \textbf{Answer all questions:} \hfill [$10 \times 2 = 20$]

\begin{enumerate}[label=\textbf{\arabic*.}]
    \item Who is the author of the book \textit{`Nationalist Thought and the Colonial World: A Derivative Discourse'}?
    \item Name the fundamental texts associated with the Indian Knowledge System (IKS) in Sociology.
    \item Define the Indian concept of \textit{Purushartha}.
    \item What is the significance of \textit{Varna} and \textit{Ashrama} in classical Indian social organization?
    \item Briefly state the core principle of \textit{Integral Humanism} (\textit{Ekatma Manavavad}).
\end{enumerate}

\section*{SECTION -- B (Short Answer Type Questions)}
\noindent \textbf{Answer any four questions:} \hfill [$4 \times 5 = 20$]

\begin{enumerate}[label=\textbf{\arabic*.}]
    \setcounter{enumii}{5}
    \item In what ways did Mahatma Gandhi's practices of fasting and simplicity reflect his interpretation of self and nation?
    \item Analyze Deendayal Upadhyay's views on the Western notion of Individualism versus the Indian concept of self (\textit{Atma}).
    \item Discuss the sociological dimensions of the Vedic perspective on environment and human relation.
    \item Explain the concept of \textit{Dharma} as a unifying social order in Indian society.
\end{enumerate}

\section*{SECTION -- C (Long Answer Type Questions)}
\noindent \textbf{Answer any two questions:} \hfill [$2 \times 15 = 30$]

\begin{enumerate}[label=\textbf{\arabic*.}]
    \setcounter{enumii}{9}
    \item Critically evaluate the relevance of `Integral Humanism' in contemporary Indian social and political thought.
    \item Discuss the key methodologies employed in Indian Knowledge Systems for understanding social reality and human behavior.
    \item Explain the evolution of Indian sociological thought from traditional Indian texts to modern perspectives.
\end{enumerate}

\end{document}
